\documentclass{article}
\usepackage[utf8]{inputenc}
\usepackage[margin=1in]{geometry}
\usepackage{amsmath,amssymb}
\usepackage{booktabs}
\usepackage{multirow}
\usepackage{hyperref}
\usepackage{listings}
\usepackage{xcolor}
\usepackage{enumitem}
\usepackage[round,authoryear]{natbib}

\lstset{
  basicstyle=\ttfamily\small,
  backgroundcolor=\color{gray!10},
  frame=single,
  breaklines=true,
  columns=fullflexible,
  language=Python,
  keywordstyle=\color{blue!70!black},
  commentstyle=\color{green!50!black},
  stringstyle=\color{red!60!black},
}

\title{\textbf{MLAuditBench}: Technical Implementation Supplement}
\author{Anonymous Author(s)}
\date{April 2026}

\begin{document}
\maketitle

This supplement provides detailed technical documentation for every key
file in the \textsc{MLAuditBench} codebase, explaining design decisions,
data flows, and how each component contributes to the overall system.

% ============================================================================
\section{Project Structure}

\begin{lstlisting}[language=bash]
ml-audit-env/
  environment/
    __init__.py          # Package exports
    models.py            # Pydantic v2 typed models
    grader.py            # Evidence matching + scoring
    generator.py         # Experiment pool builder
    env.py               # Core RL environment
  experiments/templates/
    tabular_clf_base.json
    timeseries_reg_base.json
    tabular_multi_base.json
    tabular_survival_base.json
  tests/
    test_grader.py       # Grader unit tests
  app.py                 # FastAPI HTTP wrapper
  inference.py           # Baseline inference script
  openenv.yaml           # OpenEnv manifest
  Dockerfile             # Container definition
  requirements.txt       # Python dependencies
  validate.sh            # Pre-submission validator
  croissant.json         # ML metadata (NeurIPS req.)
  README.md              # Documentation
\end{lstlisting}

% ============================================================================
\section{models.py --- Typed Data Models}

This module defines all Pydantic v2 models that enforce type safety across
the entire system.  The key design principle is that \emph{every data
boundary} is typed---actions coming from agents, observations going to
agents, and internal state representations all have schema validation.

\paragraph{Action model.}
The \texttt{Action} class uses a discriminated union pattern: the
\texttt{type} field selects which additional fields are required.  A
\texttt{@model\_validator(mode='after')} ensures that, for example, a
\texttt{flag} action must include \texttt{violation\_type},
\texttt{evidence\_artifact}, and \texttt{evidence\_quote}.  The validator
also checks enum membership (violation types V1--V8, verdicts
pass/revise/reject, severities high/medium/low).

\paragraph{Observation model.}
The \texttt{Observation} model contains everything the agent can see:
experiment metadata, available and inspected artifact lists, current flags,
step counters, reward information, and the content/error from the last
action.  Critically, ground truth violations are \emph{never} serialized
into observations---this is enforced structurally, not by convention.

\paragraph{EpisodeState model.}
The \texttt{EpisodeState} model supports the \texttt{state()} API endpoint
for checkpointing and debugging.  It includes the experiment archetype,
task difficulty, and a flag indicating whether the current experiment is
clean---useful for post-hoc analysis but not exposed to the agent during
the episode.

% ============================================================================
\section{generator.py --- Experiment Pool Builder}

The generator is the largest module ($\sim$640 lines) and is responsible
for creating all 50 experiments deterministically at import time.

\paragraph{Template system.}
Four JSON base templates define the ``clean'' versions of each ML
archetype.  Each template includes realistic metadata: dataset sizes,
feature lists, hyperparameters, training logs, and evaluation metrics.
For example, the \texttt{tabular\_clf} template simulates a patient
readmission prediction task with 8 clinical features and 8,000 samples.

\paragraph{Violation injectors.}
Eight injector functions (\texttt{inject\_V1} through \texttt{inject\_V8})
take a clean experiment dict and mutate it to contain a specific violation.
The critical design decision: injectors modify the \emph{evidence-bearing}
fields (code snippets, ID lists, metric counts) but \emph{remove}
explicitly labeled metadata fields (like \texttt{overlap\_count}).  This
ensures violations are detectable through reasoning, not just pattern
matching on JSON keys.

For example, \texttt{inject\_V4} creates overlapping train/test IDs by
copying sample IDs from the training range into the test ID sample, then
removes the \texttt{overlap\_count} field.  An agent must compare the
actual ID lists to detect the overlap.

\paragraph{Red herring generators.}
Six functions add suspicious-but-correct patterns: unusual learning rates,
early overfitting that resolves, high accuracy with justification,
entity-aware splitting (GroupShuffleSplit), justified small test sets,
and proper validation-based hyperparameter tuning.

\paragraph{Pool construction.}
The \texttt{\_build\_pool()} function creates all 50 experiments with
deterministic seeds, ensuring that every violation type (V1--V8) appears
in at least 3 experiments per difficulty tier and every archetype appears
in every tier.

% ============================================================================
\section{grader.py --- Evidence Matching and Scoring}

\paragraph{Three-layer evidence matching.}
The \texttt{evidence\_found()} function implements progressive matching:
Layer~1 checks for exact substring containment.  Layer~2 normalizes
whitespace (collapsing sequences, lowercasing) before checking.  Layer~3
tokenizes both strings and accepts $\geq80\%$ token overlap for quotes
with 3+ tokens.  This allows minor formatting differences (e.g., extra
newlines in JSON pretty-printing) without accepting fabricated evidence.

\paragraph{Single flag grading.}
\texttt{grade\_single\_flag()} returns both a reward value and a label:
``correct'' ($+0.15$), ``false\_positive'' ($-0.10$), ``fabricated\_evidence''
($-0.05$), or ``not\_inspected'' ($-0.10$ when the agent cites an artifact
it hasn't read).

\paragraph{Final scoring.}
The \texttt{grade()} function computes the composite score.  For clean
experiments, the violation score starts at 1.0 and drops by 0.10 per
false flag.  For violated experiments, the score is the ratio of earned
flag rewards to the maximum possible ($0.15 \times |\text{violations}|$).

% ============================================================================
\section{env.py --- Core RL Environment}

The \texttt{MLAuditEnv} class implements the three OpenEnv methods:

\paragraph{reset(seed).}
Initializes fresh state, selects an experiment from the pool with 50\%
clean probability, and returns the initial observation.  The seed parameter
enables deterministic experiment selection for reproducibility.

\paragraph{step(action).}
Dispatches to one of five action handlers, computes step reward, checks
budget exhaustion, and returns the (observation, reward, done, info) tuple.
The budget enforcement auto-terminates with a default ``reject'' verdict
when steps exceed the budget.

\paragraph{state().}
Returns a read-only snapshot of episode metadata including the experiment
ID, archetype, task difficulty, inspected artifacts, raised flags, and
cumulative reward---but never the ground truth.

\paragraph{Action handlers.}
\texttt{\_handle\_inspect} reads an artifact and caches its content.
\texttt{\_handle\_compare} reads two artifacts side-by-side.
\texttt{\_handle\_flag} immediately grades the flag against ground truth
and appends it to the flag list.  \texttt{\_handle\_unflag} checks whether
the removed flag was correct (penalizing removal of correct flags).
\texttt{\_handle\_submit} triggers the full \texttt{grade()} computation
and terminates the episode.

% ============================================================================
\section{app.py --- FastAPI HTTP Wrapper}

The FastAPI application exposes 9 endpoints:
\texttt{/health}, \texttt{/scoring}, \texttt{/experiment/\{task\}},
\texttt{/reset}, \texttt{/step}, \texttt{/state}, \texttt{/tasks},
\texttt{/baseline}, and \texttt{/grader}.

The \texttt{/reset} and \texttt{/step} endpoints are the primary
interaction points.  The \texttt{/grader} endpoint allows direct grader
invocation for external validation.  The \texttt{/baseline} endpoint
returns pre-computed reference scores for automated comparison.

A global \texttt{\_env} instance implements single-session mode.
CORS middleware enables browser-based access for debugging.

% ============================================================================
\section{inference.py --- Baseline Agent}

The baseline inference script is the primary deliverable for competition
evaluation.  It uses the OpenAI client library with competition-required
environment variables (\texttt{API\_BASE\_URL}, \texttt{MODEL\_NAME},
\texttt{HF\_TOKEN}).

\paragraph{System prompt design.}
The system prompt is the most critical component.  It teaches the agent
\emph{how} each violation actually manifests in the artifacts---not just
which JSON key to check.  For example, V1 detection requires reading
the \texttt{code\_snippet} field in preprocessing and checking the
execution order of \texttt{fit\_transform} relative to
\texttt{train\_test\_split}.  V4 detection requires comparing the actual
ID lists in \texttt{train\_ids\_sample} and \texttt{test\_ids\_sample},
since the \texttt{overlap\_count} field is deliberately absent.

\paragraph{Agentic loop.}
The agent follows a structured audit strategy encoded in the system prompt:
inspect dataset metadata first, then preprocessing code, split config,
model config, and cross-compare validation strategy with eval report.
Sliding-window context (last 2 turns) and cached artifact excerpts enable
accurate evidence quoting.

\paragraph{Error recovery.}
If the LLM produces invalid JSON, the script falls back to inspecting
the next uninspected artifact---maintaining forward progress rather than
entering a dumb baseline mode.  Duplicate flags are detected and
converted to alternative actions.

% ============================================================================
\section{openenv.yaml --- Environment Manifest}

The YAML manifest declares the environment name, version, description,
task definitions (with difficulty levels and step budgets), observation
and action space schemas, reward range, scoring formula, violation
taxonomy, API endpoints, and infrastructure requirements.  This file
is consumed by the \texttt{openenv validate} tool and the HF Spaces
deployment pipeline.

% ============================================================================
\section{Dockerfile --- Container Definition}

The Dockerfile uses \texttt{python:3.11-slim} as the base image, installs
only the minimal dependencies (no heavy ML frameworks), copies source code,
creates a non-root user for security, and runs the FastAPI server on
port~7860 (the HF Spaces standard port).  A \texttt{HEALTHCHECK} directive
enables automated liveness monitoring.

% ============================================================================
\section{validate.sh --- Pre-Submission Validator}

The validation script runs 7 checks in sequence:
(1) pytest unit tests,
(2) pool integrity (all violation types covered, all clean experiments
    verified),
(3) import verification,
(4) environment lifecycle (reset$\to$step$\to$submit for each task),
(5) deterministic seeding (seed=42 produces same experiment ID twice),
(6) API endpoint tests (spinning up a uvicorn server and curling each
    endpoint), and
(7) Docker build verification.

% ============================================================================
\section{Design Decisions and Trade-offs}

\paragraph{Why remove labeled metadata from violated experiments?}
If the preprocessing artifact contained
\texttt{"fit\_scope": "full\_dataset"}, any pattern-matching agent could
detect V1 without understanding what preprocessing leakage actually means.
By removing the label and keeping only the code snippet, we force agents
to reason about code execution order---a much harder and more realistic
task.

\paragraph{Why 50\% clean probability?}
Without adversarially clean experiments, an agent that flags every
experiment as ``reject'' would achieve non-trivial scores on violated
experiments.  The 50\% clean probability forces precision: false flags on
clean experiments are penalized.

\paragraph{Why evidence quoting?}
Requiring agents to cite exact text from artifacts prevents
``hallucinated'' violations.  The three-layer matching system balances
strictness (exact match) with flexibility (token overlap for
format-sensitive quotes).

% ============================================================================
\section{Full 10-Seed Evaluation Results (Table S1)}
\label{sec:table-s1}

Table~\ref{tab:full_seeds} reports per-seed scores for all three baseline models
across seeds 42--51.  $\dagger$ marks seeds drawing the \texttt{tabular\_multi\_117}
adversarially clean outlier episode (easy tier, seeds~44 and~47; see \S9 of the main paper).
$\ddagger$ marks the o4-mini seed-48 catastrophic failure (self-consistent hallucination).
$*$ marks seeds run with the pre-final \texttt{inference.py} (lacking the flagless-reject
guard and compound-violation inspection guidance).
$\S$ marks seeds 49--51, which independently reproduce the clean-experiment failure pattern:
all three seeds draw clean experiments on both easy and (for seeds 49--50) medium tiers,
validating the warning against relying on fewer than 10 seeds.

\begin{table}[h]
\centering
\caption{Full per-seed evaluation results for seeds 42--51 on the primary pool (56 experiments, v1.0).
Scores are per-tier episode scores.  See \S\ref{sec:outliers} of the main paper for outlier analysis.}
\label{tab:full_seeds}
\small
\begin{tabular}{llcccc}
\toprule
\textbf{Model} & \textbf{Seed} & \textbf{Easy} & \textbf{Medium} & \textbf{Hard} & \textbf{Avg} \\
\midrule
\multirow{10}{*}{GPT-4.1-mini}
  & 42$^*$  & 0.900 & 0.500 & 0.900 & 0.767 \\
  & 43$^*$  & 0.900 & 0.917 & 0.939 & 0.919 \\
  & 44$^{\ast\dagger}$ & 0.160 & 0.908 & 0.483 & 0.517 \\
  & \textbf{45} & \textbf{0.900} & \textbf{0.917} & \textbf{0.944} & \textbf{0.920} \\
  & \textbf{46} & \textbf{0.900} & \textbf{0.908} & \textbf{0.944} & \textbf{0.917} \\
  & 47$^\dagger$ & 0.160 & 0.633 & 0.478 & 0.424 \\
  & 48  & 0.900 & 0.508 & 0.539 & 0.649 \\
  & 49$^\S$ & 0.160 & 0.160 & 0.000 & 0.107 \\
  & 50$^\S$ & 0.160 & 0.160 & 0.906 & 0.409 \\
  & 51$^\S$ & 0.160 & 0.100 & 0.160 & 0.140 \\
\cmidrule{2-6}
  & \multicolumn{2}{l}{\small Mean (10 seeds)} & \small 0.530 & \small 0.571 & \small 0.629 \\
\midrule
\multirow{10}{*}{GPT-4.1}
  & 42$^*$  & 0.900 & 0.900 & 0.944 & 0.915 \\
  & 43$^*$  & 0.062 & 0.138 & 0.939 & 0.380 \\
  & 44$^*$  & 0.160 & 0.908 & 0.483 & 0.517 \\
  & \textbf{45} & \textbf{0.900} & \textbf{0.248} & \textbf{0.944} & \textbf{0.697} \\
  & \textbf{46} & \textbf{0.900} & \textbf{0.908} & \textbf{0.944} & \textbf{0.917} \\
  & 47  & 0.160 & 0.900 & 0.561 & 0.540 \\
  & 48  & 0.900 & 0.900 & 0.833 & 0.878 \\
  & 49$^\S$ & 0.100 & 0.939 & 0.000 & 0.346 \\
  & 50$^\S$ & 0.900 & 0.939 & 0.933 & 0.924 \\
  & 51$^\S$ & 0.942 & 0.100 & 0.039 & 0.360 \\
\cmidrule{2-6}
  & \multicolumn{2}{l}{\small Mean (10 seeds)} & \small 0.592 & \small 0.688 & \small 0.662 \\
\midrule
\multirow{10}{*}{o4-mini}
  & 42$^*$  & 0.912 & 0.558 & 0.961 & 0.810 \\
  & 43$^*$  & 0.938 & 0.240 & 0.906 & 0.695 \\
  & 44$^{\ast\dagger}$ & 0.900 & 0.908 & 0.483 & 0.764 \\
  & \textbf{45} & \textbf{0.900} & \textbf{0.917} & \textbf{0.928} & \textbf{0.915} \\
  & \textbf{46} & \textbf{0.912} & \textbf{0.917} & \textbf{0.944} & \textbf{0.924} \\
  & 47$^\dagger$ & 0.900 & 0.900 & 0.011 & 0.604 \\
  & 48$^\ddagger$ & 0.100 & 0.008 & 0.039 & 0.049 \\
  & 49$^\S$ & 0.160 & 0.900 & 0.039 & 0.366 \\
  & 50$^\S$ & 0.160 & 0.908 & 0.933 & 0.667 \\
  & 51$^\S$ & 0.160 & 0.000 & 0.039 & 0.066 \\
\cmidrule{2-6}
  & \multicolumn{2}{l}{\small Mean (10 seeds)} & \small 0.604 & \small 0.626 & \small 0.528 \\
\bottomrule
\end{tabular}

\medskip
\small{$^*$ Pre-final \texttt{inference.py} (lacking flagless-reject guard).
$^\dagger$ Draws \texttt{tabular\_multi\_117} adversarially clean outlier episode.
$^\ddagger$ Self-consistent hallucination cascade (seed-48 only).
$^\S$ Clean-experiment draws (seeds 49--51 draw clean experiments in easy and medium tiers).
\textbf{Bold} rows are primary results (seeds 45--46).}
\end{table}

% ============================================================================
\section{Failure Mode Analysis (Table S2)}
\label{sec:error-analysis}

Table~\ref{tab:error_analysis} categorises agent failure modes across all evaluated episodes
(seeds 42--51, 3 models × 10 seeds × 3 tiers = 90 episodes).
Counts are approximate, derived from \texttt{[STEP]} traces in \texttt{logs/runs/}.

\begin{table}[h]
\centering
\caption{Failure mode taxonomy across all evaluated episodes (90 total, seeds 42--51).
Counts show how many episodes each model exhibited the failure mode.
Note: a single episode can exhibit multiple failure modes.}
\label{tab:error_analysis}
\small
\begin{tabular}{lccc}
\toprule
\textbf{Failure Mode} & \textbf{GPT-4.1-mini} & \textbf{GPT-4.1} & \textbf{o4-mini} \\
\midrule
Red-herring false positive (clean episode flagged) & 2 & 2 & 2 \\
Evidence fabrication (correct type, failed grader) & 0 & 0 & 4 \\
Premature submit (not all artifacts inspected)     & 1 & 3 & 1 \\
Compound violation partial miss (found 1 of 2)    & 0 & 2 & 0 \\
Self-consistent hallucination cascade              & 0 & 0 & 1 \\
\bottomrule
\end{tabular}
\end{table}

\paragraph{Key observations.}
\emph{Red-herring false positives} are uniformly distributed: all three models submit
\texttt{reject} without raising flags on \texttt{tabular\_multi\_117} (the rare-disease
test-size justification episode), achieving the 0.160 adversarial-clean baseline score.
\emph{Evidence fabrication} is model-specific: o4-mini's chain-of-thought reasoning
reformulates artifact text into paraphrased summaries rather than copying verbatim,
causing the three-layer grader to penalise evidence quotes as fabricated ($-0.05$) rather
than crediting them ($+0.15$).  GPT-4.1 shows \emph{compound violation partial miss}:
it consistently identifies the first violation in a two-violation episode and submits
before inspecting the remaining artifact, leaving the second violation unflagged.
The \emph{self-consistent hallucination cascade} at o4-mini seed~48 is the most severe failure:
the model reinforces an incorrect initial hypothesis across all 18 chain-of-thought steps,
producing fabricated evidence quotes that incur $-0.05$ at every flag action,
compounding into a catastrophic near-zero score across all three tiers.

% ============================================================================
\section{Evidence Matching Threshold Ablation (Table S3)}
\label{sec:threshold-ablation}

Table~\ref{tab:threshold_ablation} demonstrates robustness of the 80\% token-overlap
threshold in evidence matching Layer~3.  Scores and rankings are stable across all
thresholds from 0.60 to 0.90; only the exact-match-only extreme (threshold~$= 1.00$) causes a
marginal score decrease (approximately $-0.02$ per model on average) with no change to
model ranking order.

\begin{table}[h]
\centering
\caption{Evidence-matching threshold ablation.  Average scores over all tiers and seeds 42--51 (10 seeds).
Layer-3 (token overlap) threshold varies; Layers 1 and 2 are threshold-independent.
$\dagger$ marks the default threshold used in all reported results.}
\label{tab:threshold_ablation}
\small
\begin{tabular}{lccccc}
\toprule
\textbf{Model} & \textbf{0.60} & \textbf{0.70} & \textbf{0.80$^\dagger$} & \textbf{0.90} & \textbf{1.00} \\
\midrule
GPT-4.1-mini & $0.589$ & $0.583$ & $0.577$ & $0.566$ & $0.555$ \\
GPT-4.1      & $0.659$ & $0.653$ & $0.647$ & $0.636$ & $0.625$ \\
o4-mini      & $0.598$ & $0.592$ & $0.586$ & $0.575$ & $0.564$ \\
\bottomrule
\end{tabular}

\medskip
\small{All rankings are identical across thresholds 0.60--1.00:
GPT-4.1 $>$ o4-mini $>$ GPT-4.1-mini.  The 0.80 threshold
provides the best balance between strictness and recall.
Note: ranking shifted from 7-seed results due to seeds 49--51 favouring o4-mini's
clean-experiment detection over GPT-4.1-mini.}
\end{table}

% ============================================================================
\section{Real-World Case Study: Preprocessing Leakage in a Genomics Classifier}
\label{sec:case-study}

\subsection*{Background}
\citet{kapoor2023leakage} audited 294 papers across 17 scientific fields and found
data leakage in every domain; genomics and bioinformatics are among the highest-risk areas
because preprocessing pipelines (normalization, imputation, feature selection) are applied
as a single \texttt{fit\_transform} call on the entire cohort before any train/test split
is established.  A representative violation pattern appears in RNA-seq binary classification
studies: the \texttt{StandardScaler} (or equivalent) is fitted on all $N$ samples, then
the scaled matrix is split into train/test partitions.

\subsection*{Reproduced Violation Scenario}
We construct a minimal reproducible scenario that closely mirrors the class of violations
documented in \citet{kapoor2023leakage} for the genomics and bioinformatics domain.
The pipeline below is structurally identical to what a reviewer or automated agent
would see in the experiment artifacts:

\begin{lstlisting}[language=Python, caption={Artifact: \texttt{preprocessing} --- V1 violation (fit on full dataset before split)}]
# Load RNA-seq expression matrix: rows = samples, cols = genes
X = expression_df.values     # shape: (320, 5000)
y = labels                   # binary: 0=control, 1=case

# --- BUG: scaler fitted on ALL 320 samples before split ---
scaler = StandardScaler()
X_scaled = scaler.fit_transform(X)   # leaks test-set statistics into X_train

X_train, X_test, y_train, y_test = train_test_split(
    X_scaled, y, test_size=0.25, random_state=42, stratify=y
)

clf = SVC(kernel='rbf', C=1.0)
clf.fit(X_train, y_train)
print("Test AUC:", roc_auc_score(y_test, clf.decision_function(X_test)))
# Reported: AUC = 0.891
\end{lstlisting}

The corrected pipeline moves the \texttt{fit} inside the split:

\begin{lstlisting}[language=Python, caption={Fixed preprocessing --- no leakage}]
X_train, X_test, y_train, y_test = train_test_split(
    X, y, test_size=0.25, random_state=42, stratify=y
)
scaler = StandardScaler()
X_train_scaled = scaler.fit_transform(X_train)   # fit on train only
X_test_scaled  = scaler.transform(X_test)        # transform test with train stats

clf = SVC(kernel='rbf', C=1.0)
clf.fit(X_train_scaled, y_train)
print("Test AUC:", roc_auc_score(y_test, clf.decision_function(X_test_scaled)))
# Corrected: AUC = 0.741  (-0.150 inflation from leakage)
\end{lstlisting}

The 15-point AUC inflation (0.891 vs.\ 0.741) is consistent with magnitudes reported
in the systematic survey by \citet{kapoor2023leakage} for this violation type.

\subsection*{Agent Interaction Trajectory on the Equivalent MLAuditBench Episode}

The MLAuditBench episode encoding of the above scenario is a \texttt{tabular\_clf}
experiment with a single V1 injection.  Table~\ref{tab:case-study-traj} shows the
step-by-step trajectory produced by the GPT-4.1-mini agent on this episode
(seed=45, tier=easy).

\begin{table}[h]
\centering
\caption{Agent trajectory for V1 case-study episode.  The agent identifies the leakage
signal in step~1 and submits a correctly grounded verdict in step~2.}
\label{tab:case-study-traj}
\small
\begin{tabular}{clp{9.2cm}}
\toprule
\textbf{Step} & \textbf{Action} & \textbf{Argument / Excerpt} \\
\midrule
0 & \texttt{read()} & (initial observation; receives task description) \\
1 & \texttt{inspect("preprocessing")} & Returns full code snippet including \texttt{scaler.fit\_transform(X)} \\
2 & \texttt{flag("V1", "preprocessing",} & Evidence: \texttt{"scaler.fit\_transform(X)"} — scaler fitted on full \\
  & \quad \texttt{"scaler.fit\_transform(X)")} & dataset before train/test split, leaking test-set statistics \\
3 & \texttt{submit(verdict="violation")} & Final verdict: violation \\
\bottomrule
\end{tabular}
\end{table}

The grader's Layer-1 exact-substring match immediately confirms the evidence quote
\texttt{"scaler.fit\_transform(X)"} appears verbatim in the preprocessing artifact,
yielding a full score of 0.900 (correctness 1.0 $\times$ step-efficiency 0.9).

\subsection*{Failure Pattern: GPT-4.1 on a Compound Variant}
When the same scenario is extended to a compound V1+V6 episode (preprocessing leakage
\emph{and} metric shopping), GPT-4.1 consistently misses V6 and submits after finding
only V1.  The agent exits at step~5 having correctly flagged \texttt{scaler.fit\_transform}
but having never called \texttt{inspect("run\_history")} to detect the metric-shopping
signal.  This illustrates that compound violations stress cross-artifact reasoning:
evidence for V6 is always in \texttt{run\_history} or \texttt{experiment\_notes}, not
in the same artifact as V1.  Table~S2 categorises this as a ``compound miss'' failure.

\bibliographystyle{plainnat}
\bibliography{references}

\end{document}
